\section{Conclusion}
\label{sec:conclusion}

This paper has described an open Python emulator of the OBR's macroeconometric model: 372 published EViews equations transpiled and solved in the open, anchored to the November 2025 EFO by the OBR's own add-factor device, and connected to household-resolution microsimulation through a static-costing bridge that reproduces the official ``static costing in, second-round effects out'' workflow. The anchored baseline reproduces the published forecast to 0.18 per cent on GDP and is hard-gated in CI; the free-running model is documented as weak, with every residual error traced to a fixed bug or a floor imposed by unpublished OBR inputs. A worked 1p basic-rate reform yields a costing within the range of HMRC's published ready reckoner, with small negative second-round GDP effects of the sign and size the official conventions imply.

The broader aim is methodological. Official fiscal scoring in the UK runs through a model that is described but not runnable by the public. Emulating it openly---with its judgement device made explicit, its passthroughs labelled, and its calibration failures published alongside its successes---moves the transparency frontier: researchers can now inspect the marginal propensities and error-correction speeds that shape UK costings, run counterfactual shocks through the official transmission mechanisms, and see precisely where the public information runs out. Within the PolicyEngine Macro suite, the emulator supplies the near-term, official-convention answer alongside the OLG model's long-run structural answer and the microsimulation's statutory one; scoring the same reform through all three, on a declared verification gradient, is the point.

Future work includes populating the market-sector supply block (retiring the stabilised investment closure), reviving the trade and labour-demand channels currently held exogenous, recycling income feedback through the bridge, and re-anchoring to each new EFO vintage as it is published.
